%% CV General enfocado a Java - Álvaro López Álvarez (Español)
%% CV genérico para puestos backend Java / Spring Boot. Sincronizado con main_java_en.tex.
%% Compilar con: lualatex main_java_es.tex
%% (pdflatex suele fallar en instalaciones modernas de MiKTeX con errores de fontawesome5.)

\documentclass[11pt,a4paper,sans]{moderncv}
\moderncvstyle{banking}
\moderncvcolor{blue}

% Fuerza el nombre, apellidos y los títulos de sección al azul de moderncv
% (color1). Por defecto, el estilo banking en lualatex+MiKTeX los deja en
% negro, lo cual desentona con el resto de la paleta azul.
\renewcommand*{\firstnamestyle}[1]{{\fontsize{34}{36}\bfseries\upshape\color{color1}#1}}
\renewcommand*{\lastnamestyle}[1]{{\fontsize{34}{36}\bfseries\upshape\color{color1}#1}}
\renewcommand*{\sectionstyle}[1]{{\sectionfont\color{color1}#1}}

\usepackage[utf8]{inputenc}
\usepackage{needspace}
\usepackage{hyperref}
\hypersetup{
    colorlinks=true,
    linkcolor=blue,
    filecolor=magenta,
    urlcolor=blue,
    pdftitle={Álvaro López Álvarez - CV},
    pdfpagemode=FullScreen,
}
\usepackage[scale=0.80]{geometry}
\usepackage{import}
\usepackage{eso-pic}

% datos personales
\name{Álvaro}{López Álvarez}
\address{Valencia, España}{}{}
\phone[mobile]{+34 640764553}
\email{alvaro.lopez.dev@gmail.com}
\extrainfo{\href{https://www.linkedin.com/in/alvarolopez-dev}{LinkedIn}, \href{https://github.com/alerone}{GitHub}, \href{https://alerone.github.io/portfolio}{Portfolio}}

% El estilo banking de moderncv (variante de cabecera 3) no conecta \photo -
% se ignora silenciosamente. Se superpone la foto directamente, anclada a la
% esquina de la página para que sea independiente de la cabecera centrada.
\AddToShipoutPictureFG{%
  \ifnum\value{page}=1
    \AtPageUpperLeft{%
      \hspace*{\dimexpr\paperwidth-2.3cm-1.3cm\relax}%
      \raisebox{\dimexpr-1.3cm-2.3cm\relax}{\fboxsep0pt\fboxrule0.6pt\fcolorbox{color1}{white}{\includegraphics[width=2.3cm]{photo}}}%
    }%
  \fi
}

\begin{document}

\makecvtitle

% Titular: lo que es / quiere ser el candidato, visible a primera vista.
\vspace{-4pt}
\begin{center}
{\Large\bfseries\color{color1}Desarrollador Backend Junior \textbar{} Java y Spring Boot}\\[3pt]
{\normalsize\color{color2}APIs REST \textbullet{} Microservicios Orientados a Eventos \textbullet{} Arquitectura Hexagonal}
\end{center}

% ============================================================
%     PERFIL PROFESIONAL
% ============================================================

\vspace{2pt}
\small{\textbf{Ingeniero de software especializado en el ecosistema Java}, centrado en construir APIs REST y microservicios orientados a eventos con una arquitectura limpia y testeable. Experiencia práctica con Spring Boot 3: microservicios con Kafka (patrón outbox, dead-letter topics, procesamiento idempotente), APIs REST con JPA y Jakarta Validation, y pruebas de integración con Testcontainers. Actualmente profundizando en Java mediante el estudio autodidacta de la arquitectura hexagonal (puertos y adaptadores) con Spring Boot y de patrones de diseño arquitectónico. También domino Golang, TypeScript y React, con experiencia profesional de I+D desarrollando agentes basados en LLM en el ITI. Máster en Ingeniería Informática (nota final 8,8/10) e inglés B2 certificado.}

% ============================================================
%     COMPETENCIAS CLAVE
% ============================================================

\section{Competencias Clave}
\vspace{1pt}
\begin{itemize}

\item \textbf{Backend Java}: Java, Spring Boot 3, Spring Data JPA, Jakarta Validation, APIs REST, DTOs/mappers, gestión centralizada de errores, JUnit, Testcontainers

\item \textbf{Arquitectura}: Arquitectura hexagonal (puertos y adaptadores), microservicios orientados a eventos, patrón outbox, dead-letter topics, consumidores idempotentes

\item \textbf{Mensajería e Infraestructura}: Apache Kafka, NATS, Docker, Docker Compose, Apache APISIX, PostgreSQL

\item \textbf{Otros Lenguajes y Frontend}: Golang, TypeScript/JavaScript, React, Python, Kotlin, SQL

\item \textbf{Ingeniería de IA/LLM}: LangChain, LangGraph, API de Gemini, API de OpenAI, pipelines RAG

\item \textbf{Herramientas y Prácticas}: Git, GitHub Actions, TDD, OAuth 2.0/JWT, Firebase, Supabase, Vercel, Render, NeoVim

\end{itemize}

% ============================================================
%     EXPERIENCIA PROFESIONAL
% ============================================================

\section{Experiencia Profesional}
\vspace{3pt}
\begin{itemize}

% --- Puesto más reciente ---
\needspace{5\baselineskip}
\item{\cventry{Feb 2026--Jun 2026}{Desarrollador I+D}{Instituto Tecnológico de Informática (ITI)}{Valencia, España}{}{\vspace{1pt}
\begin{itemize}
    \item Implementación de mejoras basadas en LLM para un motor de búsqueda
    \item Desarrollo de un agente multi-paso con LangGraph que convierte preguntas en lenguaje natural en consultas SQL validadas de solo lectura, para resolver políticas de acceso a datos sobre un esquema de gobernanza de datos en PostgreSQL
    \item Diseño agéntico orientado a la seguridad: selección de esquema/tablas, generación estructurada de SQL, validación de ejecución, reintento ante errores
    \item Trabajo dentro de una aplicación centrada en la gestión, gobernanza y soberanía de datos industriales
\end{itemize}}}

\vspace{3pt}

% --- Puesto anterior ---
\needspace{5\baselineskip}
\item{\cventry{Feb 2024--Jun 2024}{Desarrollador Full-Stack (Prácticas)}{Instituto Tecnológico de Informática (ITI)}{Valencia, España}{}{\vspace{1pt}
\begin{itemize}
    \item Desarrollo y mantenimiento de la aplicación corporativa de gestión de personal (Intranet)
    \item Implementación de nuevas funcionalidades de frontend con React y TypeScript, mejorando la usabilidad interna
    \item Contribución al desarrollo backend, garantizando una integración fluida entre servicios
\end{itemize}}}

\end{itemize}

% ============================================================
%     PROYECTOS JAVA DESTACADOS
% ============================================================

\section{Proyectos Java Destacados}
\vspace{1pt}
\begin{itemize}

\item \textbf{Plataforma de Venta de Entradas con Kafka y Spring Boot} -- \textit{Java, Spring Boot 3, Apache Kafka, Testcontainers}: Microservicios orientados a eventos (reserva, pago, notificación, auditoría) con patrón outbox, dead-letter topics, procesamiento idempotente y pruebas de integración con Testcontainers.

\item \textbf{English Learning API} -- \textit{Java, Spring Boot, Spring Data JPA, Jakarta Validation}: API REST para practicar phrasal verbs y collocations en inglés. Entidades JPA, DTOs y mappers, validación de peticiones, gestión centralizada de errores y retos de práctica aleatorios con seguimiento de intentos.

\item \textbf{Servidor HTTP sobre TCP (Java)} -- \textit{Java, JUnit}: Servidor HTTP/1.1 implementado directamente sobre sockets, sin frameworks, para afianzar los fundamentos de Java: flujos de E/S, concurrencia con executors y pruebas unitarias con JUnit.

\end{itemize}

% ============================================================
%     OTROS PROYECTOS
% ============================================================

\section{Otros Proyectos}
\vspace{1pt}
\begin{itemize}

\item \textbf{FaaS con Golang} -- \textit{Go, Docker, NATS JetStream, Apache APISIX}: Plataforma distribuida de Function-as-a-Service en la que workers ejecutan funciones arbitrarias empaquetadas en imágenes Docker. Colas de trabajo con NATS JetStream, gateway APISIX para autenticación JWT y limitación de tasa, y escalado horizontal con réplicas de Docker Compose.

\item \textbf{AIDucate} -- \textit{Go, Gin, PostgreSQL, React, TypeScript, API de Gemini}: Plataforma SaaS impulsada por IA que genera planes de estudio y ejercicios personalizados. Backend con autenticación JWT y PostgreSQL; frontend React en Vercel, backend en Render. \textit{(En desarrollo)}

\item \textbf{ValenMove} -- \textit{Kotlin, MVVM, JSoup}: Aplicación Android nativa que unifica el transporte público de Valencia (EMT, Metrovalencia, Valenbisi) con llegadas en tiempo real a partir de APIs públicas y datos abiertos. Trabajo de Fin de Grado, calificación 9/10.

\end{itemize}

% ============================================================
%     FORMACIÓN CONTINUA
% ============================================================

\section{Formación Continua}
\vspace{1pt}
\begin{itemize}
\item \textbf{Arquitectura Hexagonal en Java con Spring Boot} -- autoformación (en curso): puertos y adaptadores, dominio independiente del framework, adaptadores con Spring Boot.
\item \textbf{Patrones de Diseño Arquitectónico} -- autoformación (en curso): patrones de arquitectura de software y sus compromisos.
\end{itemize}

% ============================================================
%     FORMACIÓN ACADÉMICA
% ============================================================

\section{Formación Académica}
\vspace{1pt}
\begin{itemize}

\needspace{5\baselineskip}
\item{\cventry{2024--2026}{Máster Universitario en Ingeniería Informática (MUIINF)}{Universitat Politècnica de València (UPV)}{Valencia, España}{}{\vspace{1pt}
Finalizado (septiembre de 2026), nota final 8,8/10. Temario: sistemas distribuidos, computación de altas prestaciones, ciberseguridad, sistemas empotrados.
}}

\vspace{3pt}

\needspace{5\baselineskip}
\item{\cventry{2020--2024}{Grado en Ingeniería Informática, itinerario de Ingeniería del Software}{Universitat Politècnica de València (UPV)}{Valencia, España}{}{\vspace{1pt}
Nota media: 8,88/10. Trabajo de Fin de Grado (TFG): ``ValenMove -- Aplicación Unificada de Transporte Público'' (calificación 9/10).
}}

\end{itemize}

% ============================================================
%     IDIOMAS
% ============================================================

\section{Idiomas}
\vspace{1pt}
\begin{itemize}
\item Español (nativo), catalán (nativo), inglés (B2 certificado, conversación fluida).
\end{itemize}

% ============================================================
%     REFERENCIAS
% ============================================================

\section{Referencias}
\vspace{1pt}
\begin{itemize}
\item{Disponibles a petición.}
\end{itemize}

\end{document}
